\documentclass[11pt,a4paper,titlepage]{article}

% --- Language Switch: \jatrue for Japanese, \jafalse for English ---
\newif\ifja
\jatrue  % Switch to \jafalse to compile English version

% --- Page Geometry & Spacing ---
\usepackage[a4paper,margin=25mm]{geometry}
\usepackage{setspace}
\setstretch{1.15}

% --- LuaLaTeX Engine & Fonts ---
\usepackage{fontspec}
\ifja
  \usepackage{luatexja}
  \usepackage{luatexja-fontspec}
  \setmainjfont{Noto Serif CJK JP}[AutoFakeBold=1.5]
  \setsansjfont{Noto Sans CJK JP}[AutoFakeBold=1.5]
\fi
\setmainfont{Noto Serif}
\setsansfont{Noto Sans}
\setmonofont{Noto Sans Mono}[Scale=0.9]

% --- Core Packages ---
\usepackage{amsmath,amssymb}
\usepackage{booktabs}
\usepackage{tabularx}
\usepackage{listings}
\usepackage{xcolor}
\usepackage{hyperref}

% --- Syntax Highlighting / Code Block Style ---
\definecolor{codebg}{rgb}{0.96,0.97,0.98}
\definecolor{codeborder}{rgb}{0.85,0.88,0.92}
\definecolor{keywordcolor}{rgb}{0.12,0.35,0.68}
\definecolor{commentcolor}{rgb}{0.42,0.48,0.55}

\lstset{
  backgroundcolor=\color{codebg},
  basicstyle=\ttfamily\small,
  breaklines=true,
  breakatwhitespace=true,
  frame=single,
  rulecolor=\color{codeborder},
  columns=fullflexible,
  keepspaces=true,
  showstringspaces=false,
  commentstyle=\color{commentcolor},
  keywordstyle=\color{keywordcolor}\bfseries
}

% --- Hyperref Setup ---
\hypersetup{
  colorlinks=true,
  linkcolor=keywordcolor,
  urlcolor=keywordcolor,
  citecolor=keywordcolor,
  pdfborder={0 0 0}
}

% --- Document Metadata ---
\title{%
  \LARGE\textbf{%
    \ifja
      HQSS-v2: 不完全暗号文のクラウド集約と\\残余断片シャミア秘密分散（2-of-5）による極小QR耐障害性バックアッププロトコル
    \else
      HQSS-v2: Fault-Tolerant Cryptographic Key Escrow Protocol\\via Cloud Incomplete Ciphertext Aggregation and Shamir-Partitioned (2-of-5) QR Artifacts
    \fi
  }
}

\author{%
  \textbf{%
    \ifja Section 9 Inc. 研究開発技術仕様書 \else Section 9 Inc. Technical Monograph / Specification \fi
  }
}

\date{%
  \ifja 2026年10月 \else October 2026 \fi
}

\begin{document}

\maketitle

\begin{abstract}
\ifja
本モノグラフでは、暗号鍵（GnuPG/OpenPGP秘密鍵等）の超長期保管と耐障害性を極限まで高める新設計「Hybrid QR Secret Splitter v2 (HQSS-v2)」を提示する。本プロトコルは、共通鍵暗号化（AES-256）を施した秘密鍵を10個の等長ブロックに分割し、先頭90\%（9ブロック）を単一のバンドルとして複数のパブリッククラウドに分散配備する。欠落させた残余の10\%（1ブロック）に対してのみガロア体 $GF(2^8)$ 上のShamirの秘密分散法（2-of-5閾値）を適用し、極小の低密度QRコードとして物理媒体に分散保管する。本方式では、クラウド上の90\%を入手した攻撃者に対しても復号検証オラクルを完全に奪い、情報理論的安全性を提供する。また、物理媒体の1枚の盗難に耐えつつ、5枚中任意の2枚のスキャンとクラウド取得のみで、パスフレーズ記憶負荷を最小化しながら確実な自律復旧を可能とする。
\else
This monograph presents the Hybrid QR Secret Splitter v2 (HQSS-v2), a resilient cryptographic key escrow protocol optimized for disaster recovery and minimal cognitive overhead. HQSS-v2 fragments an AES-256 encrypted master key into ten deterministic chunks. The initial 90\% (chunks 1--9) is concatenated into a public cloud bundle and mirrored across distributed object storage. The vital remaining 10\% (chunk 10) is partitioned using Shamir's Secret Sharing over $GF(2^8)$ under a 2-of-5 threshold scheme, subsequently rendered into low-density physical QR tokens. Even in the event of an adversary harvesting the entire 90\% public bundle, the absence of padding, authentication tags, and structural metadata destroys any decryption oracle, reducing cipher analysis to information-theoretic infeasibility. Reconstruction requires merely any 2 of the 5 QR artifacts to restore the missing chunk and concatenate the local master key without third-party reliance.
\fi
\end{abstract}

\vspace{1.5em}

\section{%
  \ifja 1. 設計思想と問題提起 \else 1. Problem Statement and Design Philosophy \fi
}
\ifja
暗号化秘密鍵のバックアップは、常に「クラウド事業者の侵害・検閲リスク」と「物理メディア（紙や金属板）の破損・紛失・容量制限」の板挟みとなってきた。また、鍵長が大きく証明書チェーンを内包するGnuPG鍵をShamirの秘密分散法（SSS）でそのまま物理QRコード化する場合、バージョン（セル密度）が高くなりすぎて印刷潰れやカメラのスキャンミスを引き起こす。さらに、人間が複雑な高エントロピーパスワードを長期間記憶し続けること自体が重大な単一障害点（SPOF）となる。

HQSS-v2は、以下の3つの原則を融合してこの問題を根本解決する。
\begin{enumerate}
  \item \textbf{不完全暗号文によるオラクル消滅:} 暗号化バイナリの先頭90\%のみをクラウドに預ける。検証可能なパディングや末尾パケットを欠損させることで、攻撃者が総当たり（ブルートフォース）を試みるための正否判定オラクルを物理的に消滅させる。これにより、暗号化パスフレーズの複雑性への依存を大幅に緩和する。
  \item \textbf{質量のクラウド委任とエントロピーの物理逃がし:} 全体データの大半（90\%）をクラウドに逃がすことで、物理媒体に書き込む残余データサイズを数百バイトに抑え、低密度（Version 5--10程度）で高耐久なQRコードを実現する。
  \item \textbf{2-of-5 閾値分散による耐奪取と耐喪失の両立:} 手元に残す重要ブロックのみに $GF(2^8)$ シャミア秘密分散を適用する。1枚のQRコードの盗難・押収では情報漏洩が一切起きず、かつ5枚中3枚が焼失・紛失しても2枚から完全復旧を可能にする。
\end{enumerate}
\else
Key escrow architectures consistently clash against two failure modes: cloud provider compromise/censorship, and physical media degradation or volumetric limits. Encoding an entire GnuPG secret keyring via Shamir's Secret Sharing (SSS) directly onto paper yields high-density QR versions that suffer catastrophic read-failure rates under harsh optical conditions. Furthermore, relying on human cognitive retention of high-entropy passphrases introduces a fatal single point of failure (SPOF).

HQSS-v2 resolves these contradictions through three core design axioms:
\begin{enumerate}
  \item \textbf{Oracle Annihilation via Incomplete Ciphertext:} Releasing only 90\% of the ciphertext to untrusted clouds strips away MAC tags, padding, and packet trailers, entirely denying adversaries a mathematical verification oracle for brute-force attacks.
  \item \textbf{Bulk Offloading vs. Pure Entropy Escrow:} Offloading the vast majority of volume to public infrastructure allows physical shares to remain under a few hundred bytes, producing robust, low-version, fault-tolerant QR matrices.
  \item \textbf{2-of-5 Threshold Duality:} Confining the Shamir scheme solely to the remaining 10\% fragment delivers mathematical resistance against single-token seizure while surviving the loss of up to 3 physical cards.
\end{enumerate}
\fi

\section{%
  \ifja 2. プロトコル仕様とパイプライン \else 2. Protocol Specification and Pipeline \fi
}

\subsection{%
  \ifja 2.1 暗号化・分割・物理トークン生成 \else 2.1 Encryption, Fragmentation, and Token Generation \fi
}
\ifja
マスター秘密鍵平文 $M$ から配布用アーティファクトを生成する手続きは以下の通りである。
\begin{enumerate}
  \item \textbf{暗号化:} パスフレーズ $P$（またはランダム生成鍵）を用い、AES-256（CBCまたはGCM）およびArmored形式で $M$ を暗号化し、暗号文 $C$ を生成する。
  \item \textbf{10分割スライス:} GNU \texttt{split -n 10} 等を用い、$C$ をバイト均等に10個の断片 $\{B_0, B_1, \dots, B_9\}$ に分割する。
  \item \textbf{クラウドバンドルの生成:} 先頭9個の断片を結合し、クラウド用バルクデータ $B_{\text{cloud}} = B_0 \mathbin{\Vert} \dots \mathbin{\Vert} B_8$ を作成する。これを複数の外部オブジェクトストレージ（Google Drive, Nextcloud, S3等）に同一内容で配置する。
  \item \textbf{残余10\%のシャミア秘密分散:} 最後の1ブロック $B_9$ に対してのみ、ガロア体 $GF(2^8)$ 上で $(k, n) = (2, 5)$ のShamir秘密分散を実行し、5つのシェア $\{S_1, S_2, S_3, S_4, S_5\}$ を生成する。
  \item \textbf{QRコード印字:} 各シェア $S_i$ をBase64文字列化し、\texttt{qrencode} を用いて中解像度QRコード（誤り訂正レベルM推奨）として出力・印刷する。
\end{enumerate}
\else
The generation pipeline executes deterministically as follows:
\begin{enumerate}
  \item \textbf{Symmetric Encryption:} Encrypt plaintext secret key $M$ using AES-256 with passphrase $P$ to produce armored ciphertext $C$.
  \item \textbf{Decary Fragmentation:} Slice $C$ evenly into 10 byte segments $\{B_0, B_1, \dots, B_9\}$ via \texttt{split -n 10}.
  \item \textbf{Cloud Aggregation:} Concatenate the first 9 chunks into a single bulk artifact $B_{\text{cloud}} = B_0 \mathbin{\Vert} \dots \mathbin{\Vert} B_8$. Mirror $B_{\text{cloud}}$ across disparate public storage providers.
  \item \textbf{Shamir Dispersion of Remainder:} Subject the final chunk $B_9$ to Shamir's Secret Sharing over $GF(2^8)$ parameterized at $(k=2, n=5)$, obtaining shares $\{S_1, \dots, S_5\}$.
  \item \textbf{QR Artifact Synthesis:} Encode each share $S_i$ into Base64 format and render via \texttt{qrencode} (Error Correction Level M).
\end{enumerate}
\fi

\begin{lstlisting}[language=bash, caption={Reference Implementation: Fragmentation and QR Encoding}]
# 1. 秘密鍵のエクスポートとAES暗号化
gpg --armor --export-secret-keys <KEY-ID> | \
  gpg --symmetric --cipher-algo AES256 --armor --output key.asc

# 2. 10等分スライス
mkdir -p /tmp/hqss && cd /tmp/hqss
split -n 10 -d key.asc part_

# 3. 先頭90%を結合してクラウドバンドルを作成
cat part_00 part_01 part_02 part_03 part_04 part_05 part_06 part_07 part_08 > cloud_bundle.bin

# 4. 残余10%（part_09）を 2-of-5 シャミア秘密分散
gfsplit -n 5 -m 2 part_09 share_

# 5. Base64化してQRコード（PNG）を生成
for f in share_*; do
    base64 -w 0 "$f" | qrencode -s 6 -l M -o "${f}.png"
done
\end{lstlisting}

\subsection{%
  \ifja 2.2 復元フェーズ \else 2.2 Reconstruction Phase \fi
}
\ifja
被災時または別環境でのリカバリは、最小限の操作で完結する。
\begin{enumerate}
  \item クラウドストレージから $B_{\text{cloud}}$ を取得する。
  \item 分散保管された5枚のQRコードのうち、手元にある任意の2枚を光学スキャンし、Base64文字列からバイナリシェア $S_x, S_y$ を復元する。
  \item \texttt{gfcombine} を実行し、残余ブロック $B_9$ を数学的に復元する。
  \item $B_{\text{cloud}}$ と $B_9$ をバイナリ結合して元の暗号文 $C$ を復元し、パスフレーズ $P$ を入力して秘密鍵平文 $M$ を取り出す。
\end{enumerate}
\else
Reconstruction operates deterministically in an offline, air-gapped terminal:
\begin{enumerate}
  \item Retrieve $B_{\text{cloud}}$ from any reachable cloud mirror.
  \item Scan any 2 available physical QR cards from the 5 issued, decoding Base64 strings to yield shares $S_x, S_y$.
  \item Execute \texttt{gfcombine} with $S_x, S_y$ to mathematically restore $B_9$.
  \item Concatenate $B_{\text{cloud}} \mathbin{\Vert} B_9$ to reform ciphertext $C$, subsequently decrypting with passphrase $P$ to obtain $M$.
\end{enumerate}
\fi

\begin{lstlisting}[language=bash, caption={Reference Implementation: Reconstruction Pipeline}]
# 1. スキャンした2片をデコード
echo "$SCAN_A" | base64 -d > share_rec_1
echo "$SCAN_B" | base64 -d > share_rec_2

# 2. 2-of-5 復元による第10ブロック再生
gfcombine -o restored_part_09 share_rec_1 share_rec_2

# 3. クラウドバンドルと結合
cat cloud_bundle.bin restored_part_09 > restored_key.asc

# 4. 復号とGPGインポート
gpg --decrypt restored_key.asc | gpg --import
\end{lstlisting}

\section{%
  \ifja 3. 暗号学的安全性と脅威モデル \else 3. Cryptographic Security and Threat Model \fi
}

\begin{table}[htbp]
\centering
\caption{%
  \ifja 脅威モデルと対抗措置（HQSS-v2） \else Threat Vectors and Mitigations (HQSS-v2) \fi
}
\vspace{0.5em}
\begin{tabularx}{\textwidth}{lXX}
\toprule
\textbf{\ifja 攻撃ベクトル \else Threat Vector \fi} & 
\textbf{\ifja 影響と攻撃者の試み \else Impact / Adversary Capability \fi} & 
\textbf{\ifja 本プロトコルの防衛機構 \else Defensive Guarantee \fi} \\
\midrule
\ifja クラウド完全侵害 \else Cloud Infrastructure Seizure \fi & 
\ifja クラウド上の $B_{\text{cloud}}$（90\%）が敵対者に完全取得される。 \else Adversary intercepts full 90\% bulk ciphertext bundle. \fi & 
\ifja ファイル構造および末尾パディングが欠落しており、復号正否判定オラクルが消滅。総当たりは数学的に不成立。 \else No plaintext oracle exists; missing final block invalidates brute-force verification over $P$. \fi \\
\addlinespace
\ifja 物理QRコードの単一奪取 \else Single Physical QR Theft \fi & 
\ifja 捜索や空き巣等によりQRコードが1枚物理的に盗まれる。 \else Adversary acquires a single printed QR code share. \fi & 
\ifja 2-of-5 シャミア秘密分散の性質上、1片から得られる $B_9$ の情報量は数学的に0ビット（情報理論的安全性）。 \else Information-theoretic security: single share reveals 0 bits of information regarding $B_9$. \fi \\
\addlinespace
\ifja 災害・火災・紛失 \else Physical Degradation / Disaster \fi & 
\ifja 自宅被災等により保管していたQRコードの大半（最大3枚）を焼失。 \else Fire or environmental damage destroys up to 3 physical cards. \fi & 
\ifja 任意の2枚さえ残存していれば、多項式補間により100\%元の暗号文を復元可能（高い耐障害性）。 \else Any 2 valid shares mathematically reconstruct $B_9$ through Lagrange interpolation. \fi \\
\addlinespace
\ifja パスフレーズ忘却・認知低下 \else Human Cognitive Degradation \fi & 
\ifja 本人が複雑なパスワードを記憶できない、または忘却する。 \else Operator forgets complex passphrase. \fi & 
\ifja 暗号強度をデータ欠損が保証するため、パスフレーズは単純なPINや記憶可能な語句に簡略化可能。 \else Security relies on missing physical shares rather than high-entropy password memorization. \fi \\
\bottomrule
\end{tabularx}
\end{table}

\section{%
  \ifja 4. 考察と運用上の推奨事項 \else 4. Discussion and Operational Recommendations \fi
}
\ifja
\subsection{完全性検証メタデータ}
単純なバイナリ結合を行うため、復元時に1バイトでもオフセットや順序の狂いが生じると復号に失敗する。
この対策として、暗号化完了時のオリジナルハッシュ（\texttt{sha256sum key.asc} の先頭8--16進数文字程度）を、各QRコードの余白またはBase64ヘッダに平文メタデータとして併記しておく運用が強く推奨される。

\subsection{クラウドデータのビットロット対策}
パブリッククラウド上の $B_{\text{cloud}}$ の長期的なビット反転（Silent Data Corruption）やストレージ障害を想定し、Par2（Parchive）等のリード・ソロモン消去符号による数\%の修復ブロックを生成してクラウド上に併置しておくことで、ネットワーク側の可用性をほぼ100\%に引き上げることが可能である。
\else
\subsection{Integrity Verification Metadata}
Because restoration relies on byte-exact concatenation, any transposition error invalidates decryption. Operators are strongly advised to print the truncated SHA-256 digest of the original \texttt{key.asc} in human-readable plain text alongside the QR codes to permit rapid sanity checks prior to invoking decryption.

\subsection{Mitigating Cloud Bit-Rot}
To guard against silent data corruption across remote endpoints, generating a minor parity volume via Par2 (Reed-Solomon erasure coding) alongside $B_{\text{cloud}}$ elevates remote repository availability to near certainty without disclosing cryptographic structure.
\fi

\section{%
  \ifja 5. 結論 \else 5. Conclusion \fi
}
\ifja
HQSS-v2は、「計算量的安全性」に過度に依存した従来のパスワード暗号化管理から脱却し、「物理的な不完全暗号文によるオラクル消滅」と「ガロア体上のShamir秘密分散」を組み合わせることで、極めて高次元の耐奪取性と耐喪失性を両立した。
記憶の認知的負荷を極小化しつつ、クラウドの容量メリットと紙媒体の非通電オフライン保管メリットを完璧に抽出した、恒久的な秘密鍵バックアップの決定解である。
\else
HQSS-v2 breaks the traditional reliance on high-entropy passphrase memorization by bridging incomplete ciphertext fragmentation with Galois-field Shamir secret sharing. By removing decryption verification oracles from untrusted cloud stores and compressing the physical escrow requirement to low-density 2-of-5 QR shares, HQSS-v2 achieves an unprecedented balance between seizure-resistance, disaster recovery availability, and zero-cognitive overhead.
\fi

\end{document}
